\section{Operator Implementation}
\label{sec:operator}

\subsection{Role and Lifecycle}

Each Operator is a long-running Go binary, derived from the
Incredible Squaring template and registered with the
\texttt{PriceOracleServiceManager} at startup.  It subscribes to the
\texttt{NewTaskCreated} event stream from the TaskManager.  On each
event it parses the requested \texttt{assetPair}, queries a single
statically configured CEX endpoint, signs the resulting
\texttt{TaskResponse} with its EigenLayer-registered BLS private key,
and forwards the signature over JSON-RPC to the Aggregator's
\texttt{ProcessSignedTaskResponse} handler.  Per-Operator state is
deliberately minimal: no persistent task history is required, so
restarted Operators rejoin without replaying past tasks.

\subsection{Heterogeneous CEX Integration}

Our deployment runs four Operators, each pinned to a different
exchange via the \texttt{OPERATOR\_CEX\_SOURCE} environment variable:
Binance, Coinbase, Kraken, and OKX.  All four expose unauthenticated
ticker endpoints returning JSON.  Listing~\ref{lst:fetch} shows the
Binance fetcher; the other three are structurally similar with
exchange-specific path and field-name differences.  Errors (HTTP
non-2xx, JSON parse failure, fixed-point overflow) abort that
Operator's task without poisoning the aggregator: the missing
report simply lowers the count of available quotes for the round.

\begin{lstlisting}[language=Go, caption={CEX price fetcher (Binance).}, label={lst:fetch}]
func fetchBinance(p string) (*big.Int, error) {
  u := "https://api.binance.com/api/v3/ticker/price?symbol=" + norm(p)
  ctx, c := context.WithTimeout(context.Background(), 2*time.Second)
  defer c()
  req, _ := http.NewRequestWithContext(ctx, "GET", u, nil)
  resp, err := http.DefaultClient.Do(req)
  if err != nil { return nil, err }
  defer resp.Body.Close()
  var b struct{ Price string }
  json.NewDecoder(resp.Body).Decode(&b)
  return parseSixDecimal(b.Price)
}
\end{lstlisting}

\subsection{BLS Signing and RPC Submission}

After parsing the ticker, the Operator constructs a
\texttt{TaskResponse}, ABI-encodes it, and signs the digest with its
BLS12-381 private key using the EigenLayer SDK's BLS package.  The
signature is forwarded alongside the response and the registered
\texttt{OperatorId} to the Aggregator over JSON-RPC.  The AVS
contract is the canonical source of truth for the registered key
set, so no additional public-key plumbing is required on the
Operator side.

\subsection{Decentralization Beyond Restaking}

EigenLayer restaking provides \emph{economic} decentralization: any
holder of restaked ETH can plausibly run an Operator.  Our deployment
adds a second axis, \emph{reporter diversity}, by sourcing each
Operator's quote from a different exchange.  A single CEX going down
(maintenance, rate-limit ban, regional block) reduces the available
quote count by at most $25\%$ but cannot, on its own, push the
consensus median outside the slashing band.  This is the property
formalised by the breakdown-point analysis of
Section~\ref{sec:agg-median}.

\subsection{Measured CEX Latency}

Figure~\ref{fig:cex-latency} reports the per-Operator quote latency
distribution observed over $1{,}000$ sequential ticker calls per
exchange.  Median latency is \textit{TBD ms} for the fastest provider
and \textit{TBD ms} for the slowest; the spread is dominated by the
geographic distance to each exchange's edge servers rather than JSON
parsing or signing overhead, which together account for less than
$10\,\mathrm{ms}$.

\begin{figure}[t]
  \centering
  \framebox[0.95\columnwidth][c]{\raisebox{1.2cm}{\textit{[placeholder: per-CEX latency boxplot --- M2 to add]}}}
  \caption{End-to-end latency distribution per exchange API,
           measured over $1{,}000$ sequential calls per source.
           Honest disagreement among feeds is a few basis points,
           well below our $\tau{=}5\%$ slashing tolerance.}
  \label{fig:cex-latency}
\end{figure}
